\documentclass[10pt, letterpaper]{article}

% Packages:
\usepackage[
    ignoreheadfoot, % set margins without considering header and footer
    top=2 cm, % seperation between body and page edge from the top
    bottom=2 cm, % seperation between body and page edge from the bottom
    left=2 cm, % seperation between body and page edge from the left
    right=2 cm, % seperation between body and page edge from the right
    footskip=1.0 cm, % seperation between body and footer
    % showframe % for debugging
]{geometry} % for adjusting page geometry
\usepackage{titlesec} % for customizing section titles
\usepackage{tabularx} % for making tables with fixed width columns
\usepackage{array} % tabularx requires this
\usepackage[dvipsnames]{xcolor} % for coloring text
\definecolor{primaryColor}{RGB}{0, 0, 0} % define primary color
\usepackage{enumitem} % for customizing lists
\usepackage{fontawesome5} % for using icons
\usepackage{amsmath} % for math
\usepackage[
    pdftitle={Juan Manuel Copia's CV},
    pdfauthor={Juan Manuel Copia},
    pdfcreator={LaTeX with RenderCV},
    colorlinks=true,
    urlcolor=primaryColor
]{hyperref} % for links, metadata and bookmarks
\usepackage[pscoord]{eso-pic} % for floating text on the page
\usepackage{calc} % for calculating lengths
\usepackage{bookmark} % for bookmarks
\usepackage{lastpage} % for getting the total number of pages
\usepackage{changepage} % for one column entries (adjustwidth environment)
\usepackage{paracol} % for two and three column entries
\usepackage{ifthen} % for conditional statements
\usepackage{needspace} % for avoiding page brake right after the section title
\usepackage{iftex} % check if engine is pdflatex, xetex or luatex

% Ensure that generate pdf is machine readable/ATS parsable:
\ifPDFTeX
    \input{glyphtounicode}
    \pdfgentounicode=1
    \usepackage[T1]{fontenc}
    \usepackage[utf8]{inputenc}
    \usepackage{lmodern}
\fi

\usepackage{charter}

% Some settings:
\raggedright
\AtBeginEnvironment{adjustwidth}{\partopsep0pt} % remove space before adjustwidth environment
\pagestyle{empty} % no header or footer
\setcounter{secnumdepth}{0} % no section numbering
\setlength{\parindent}{0pt} % no indentation
\setlength{\topskip}{0pt} % no top skip
\setlength{\columnsep}{0.15cm} % set column seperation
\pagenumbering{gobble} % no page numbering

\titleformat{\section}{\needspace{4\baselineskip}\bfseries\large}{}{0pt}{}[\vspace{1pt}\titlerule]

\titlespacing{\section}{
    % left space:
    -1pt
}{
    % top space:
    0.3 cm
}{
    % bottom space:
    0.2 cm
} % section title spacing

\renewcommand\labelitemi{$\vcenter{\hbox{\small$\bullet$}}$} % custom bullet points
\newenvironment{highlights}{
    \begin{itemize}[
        topsep=0.10 cm,
        parsep=0.10 cm,
        partopsep=0pt,
        itemsep=0pt,
        leftmargin=0 cm + 10pt
    ]
}{
    \end{itemize}
} % new environment for highlights


\newenvironment{highlightsforbulletentries}{
    \begin{itemize}[
        topsep=0.10 cm,
        parsep=0.10 cm,
        partopsep=0pt,
        itemsep=0pt,
        leftmargin=10pt
    ]
}{
    \end{itemize}
} % new environment for highlights for bullet entries

\newenvironment{onecolentry}{
    \begin{adjustwidth}{
        0 cm + 0.00001 cm
    }{
        0 cm + 0.00001 cm
    }
}{
    \end{adjustwidth}
} % new environment for one column entries

\newenvironment{twocolentry}[2][]{
    \onecolentry
    \def\secondColumn{#2}
    \setcolumnwidth{\fill, 4.5 cm}
    \begin{paracol}{2}
}{
    \switchcolumn \raggedleft \secondColumn
    \end{paracol}
    \endonecolentry
} % new environment for two column entries

\newenvironment{threecolentry}[3][]{
    \onecolentry
    \def\thirdColumn{#3}
    \setcolumnwidth{, \fill, 4.5 cm}
    \begin{paracol}{3}
    {\raggedright #2} \switchcolumn
}{
    \switchcolumn \raggedleft \thirdColumn
    \end{paracol}
    \endonecolentry
} % new environment for three column entries

\newenvironment{header}{
    \setlength{\topsep}{0pt}\par\kern\topsep\centering\linespread{1.5}
}{
    \par\kern\topsep
} % new environment for the header

\newcommand{\placelastupdatedtext}{% \placetextbox{<horizontal pos>}{<vertical pos>}{<stuff>}
  \AddToShipoutPictureFG*{% Add <stuff> to current page foreground
    \put(
        \LenToUnit{\paperwidth-2 cm-0 cm+0.05cm},
        \LenToUnit{\paperheight-1.0 cm}
    ){\vtop{{\null}\makebox[0pt][c]{
        \small\color{gray}\textit{Last updated in September 2024}\hspace{\widthof{Last updated in September 2024}}
    }}}%
  }%
}%

% save the original href command in a new command:
\let\hrefWithoutArrow\href

% new command for external links:


\begin{document}
    \newcommand{\AND}{\unskip
        \cleaders\copy\ANDbox\hskip\wd\ANDbox
        \ignorespaces
    }
    \newsavebox\ANDbox
    \sbox\ANDbox{$|$}

    \begin{header}
        \fontsize{25 pt}{25 pt}\selectfont Juan Manuel Copia

        \vspace{4 pt}

        \normalsize
        \mbox{Madrid, Spain}%
        \kern 5.0 pt%
        \AND%
        \kern 5.0 pt%
        \mbox{\hrefWithoutArrow{mailto:jmcopia96@gmail.com}{jmcopia96@gmail.com}}%
        \kern 5.0 pt%
        \AND%
        \kern 5.0 pt%
        \mbox{\hrefWithoutArrow{https://juanmacopia.github.io/}{juanmacopia.github.io}}%
        \kern 5.0 pt%
        \AND%
        \kern 5.0 pt%
        \mbox{\hrefWithoutArrow{https://www.linkedin.com/in/juancopia}{linkedin.com/in/juancopia}}%
        \kern 5.0 pt%
        \AND%
        \kern 5.0 pt%
        \mbox{\hrefWithoutArrow{https://github.com/JuanmaCopia}{github.com/JuanmaCopia}}%
    \end{header}

    \vspace{5 pt - 0.4 cm}


% =================================================================================================
% -------------------------------------------------------------------------------------------------
% -----------------------------------=====>   SUMMARY   <=====-------------------------------------
% -------------------------------------------------------------------------------------------------
% =================================================================================================


    \section{Summary}

    \begin{onecolentry}
    \textit{Software engineer and researcher completing a PhD in Computer Science at Universidad Politécnica de Madrid, with hands-on experience at Amazon. I enjoy connecting academic innovation with real-world engineering—combining research in program analysis and automated reasoning with the design and deployment of scalable cloud systems. I’m passionate about the future of artificial intelligence and about building agentic software that’s autonomous, adaptive, and impactful.}
    \end{onecolentry}



% =================================================================================================
% -------------------------------------------------------------------------------------------------
% ----------------------------------=====>   EXPERIENCE   <=====-----------------------------------
% -------------------------------------------------------------------------------------------------
% =================================================================================================

    \vspace{0.1 cm}

    \section{Experience}


        \vspace{0.1 cm}
        

        \begin{twocolentry}{
            May 2025 – Present
        }
            \textbf{Software Development Engineer}, \textit{Amazon -- Madrid, Spain} \end{twocolentry}

        \vspace{0.10 cm}
        \begin{onecolentry}
            \begin{highlights}
                \item Fast ramp-up into Amazon’s large-scale development environment, mastering internal tools, deployment processes, and operational best practices.
                \item Designed, implemented, and deployed high-impact backend services using \textbf{AWS technologies} such as \textbf{EC2}, \textbf{Lambda}, \textbf{S3}, \textbf{DynamoDB}, and \textbf{CloudWatch}.
                \item Developed and maintained \textbf{infrastructure as code} using \textbf{AWS CDK (TypeScript)}, ensuring reproducibility and scalability across environments.
                \item Automated complex workflows using AI agents and agentic systems
                \item Contributed to \textbf{continuous integration and deployment pipelines}, improving reliability and reducing release cycle time.
                \item Participated in \textbf{peer code reviews}, ensuring code quality, maintainability, and adherence to Amazon’s software engineering standards.
                \item Joined the team’s \textbf{on-call rotation}, providing 24/7 availability to resolve service incidents and maintain system reliability.
                \item Mentored new hires as an onboarding buddy, helping them integrate into the team and learn Amazon’s development and operational practices.
            \end{highlights}
        \end{onecolentry}
        
        
        \vspace{0.2 cm}


        \begin{twocolentry}{
            Dec 2020 – April 2025
        }
            \textbf{Research Assistant}, \textit{IMDEA Software Institute -- Madrid, Spain} \end{twocolentry}

        \vspace{0.10 cm}
        \begin{onecolentry}
            \begin{highlights}
                \item Designed and implemented a novel Symbolic Execution
                framework to handle complex heap-allocated inputs. Addressed a
                key challenge in lazy initialization, improving analysis
                precision and eliminating false alarms.
                \item Developed an efficient solver of structural constraints
                of partially symbolic heaps.
                \item Developed and applied search-based techniques to automate precondition inference.
                \item Developed an AI-driven tool that combines Large Language Models (LLMs), random test input generation, and program analysis to improve patch correctness validation. This work integrates ML-based reasoning with automated software testing.
            \end{highlights}
        \end{onecolentry}



        \vspace{0.2 cm}



        \begin{twocolentry}{
            Oct 2019 – Oct 2020
        }
            \textbf{Research Scholarship}, \textit{Department of Computer Science -- UNRC, Argentina} \end{twocolentry}

        \vspace{0.10 cm}
        \begin{onecolentry}
            \begin{highlights}
                \item Built a Python-based Symbolic Execution engine with
                support for lazy initialization, enabling automated reasoning
                for heap-allocated objects.
                \item Applied the Symbolic Execution engine to generate
                automated test inputs, successfully identifying a critical bug
                in an educational open-source data structure.
            \end{highlights}
        \end{onecolentry}



        \vspace{0.2 cm}



        \begin{twocolentry}{
            Jan 2019 – Mar 2019
        }
            \textbf{Summer Internship}, \textit{McAfee -- Córdoba, Argentina} \end{twocolentry}

        \vspace{0.10 cm}
        \begin{onecolentry}
            \begin{highlights}
                \item Performed fuzz testing on internal software components,
                uncovering crashes and weaknesses that enhanced the software's
                robustness and reliability.
            \end{highlights}
        \end{onecolentry}



        \vspace{0.2 cm}



        \begin{twocolentry}{
            Aug 2018 – Aug 2019
        }
            \textbf{Student Teaching Assistant}, \textit{Department of Computer Science -- UNRC, Argentina} \end{twocolentry}

        \vspace{0.10 cm}
        \begin{onecolentry}
            \begin{highlights}
                \item Supported undergraduate students in 'Data Structures' and
                'Algorithms', focusing on algorithm design, implementation, and
                problem-solving.
            \end{highlights}
        \end{onecolentry}


% =================================================================================================
% -------------------------------------------------------------------------------------------------
% ----------------------------------=====>   EDUCATION   <=====------------------------------------
% -------------------------------------------------------------------------------------------------
% =================================================================================================

    \vspace{0.1 cm}

    \section{Education}

        \vspace{0.1 cm}

        \begin{twocolentry}{
            Oct 2021 – Present
        }
            \textbf{PhD in Computer Science,} \textit{Universidad Politécnica de Madrid -- Madrid, Spain}
        \end{twocolentry}

        \vspace{0.2 cm}

        \begin{twocolentry}{
            Mar 2015 – Dec 2020
        }
            \textbf{Undergraduate degree in Computer Science (5 years +
            thesis)}
        \end{twocolentry}

        \begin{onecolentry}
            \textit{Department of Computer Science, UNRC, Argentina}
            \begin{highlights}
                \item GPA: 9.02/10.0
            \end{highlights}
        \end{onecolentry}

        \vspace{0.2 cm}

        \begin{twocolentry}{
            Mar 2015 – Jun 2018
        }
            \textbf{Undergraduate degree in Computer Science (3 years +
            final project)}
        \end{twocolentry}

        \begin{onecolentry}
            \textit{Department of Computer Science, UNRC, Argentina}
            \begin{highlights}
                \item GPA: 8.81/10.0
            \end{highlights}
        \end{onecolentry}


% =================================================================================================
% -------------------------------------------------------------------------------------------------
% ---------------------------------=====>   PUBLICATIONS   <=====----------------------------------
% -------------------------------------------------------------------------------------------------
% =================================================================================================

\vspace{0.1 cm}

    \section{Publications}

        \vspace{0.1 cm}

        \begin{samepage}
            \begin{twocolentry}{
                Coming in Nov 2025
            }
                \textbf{Search-based Inference of Class Invariants: How far can Simulated Annealing take us?}
            \end{twocolentry}

            \vspace{0.1 cm}
            \begin{onecolentry}
                \mbox{\textbf{\textit{J. M. Copia}}}, \mbox{F. Molina}, \mbox{N.
                Aguirre}, \mbox{M. Frias}, \mbox{A. Gorla}, \mbox{P. Ponzio}. \\
                \textit{To appear in Symposium on Search-Based Software Engineering (SSBSE) 2025, Seoul, Republic of Korea}

            \end{onecolentry}

        \end{samepage}

        \vspace{0.2 cm}

        \begin{samepage}
            \begin{twocolentry}{
                Aug 2025
            }
                \textbf{Search-based Inference of Class Invariants (Poster)}
            \end{twocolentry}

            \vspace{0.1 cm}
            \begin{onecolentry}
                \mbox{\textbf{\textit{J. M. Copia}}}, \mbox{F. Molina}, \mbox{N.
                Aguirre}, \mbox{M. Frias}, \mbox{A. Gorla}, \mbox{P. Ponzio}. \\
                \textit{2025 Genetic and Evolutionary Computation Conference (GECCO), Malaga, Spain} \\ \textit{doi:}
                \href{https://doi.org/10.1145/3712255.3726698}{\textit{10.1145/3712255.3726698}}

            \end{onecolentry}

        \end{samepage}

        \vspace{0.2 cm}

        \begin{samepage}
            \begin{twocolentry}{
                May 2024
            }
                \textbf{Improving Patch Correctness Analysis via Random Testing and Large
                Language Models}
            \end{twocolentry}

            \vspace{0.1 cm}
            \begin{onecolentry}
                \mbox{F. Molina}, \mbox{\textbf{\textit{J. M. Copia}}}, \mbox{A.
                Gorla}. \\ \textit{2024 IEEE Conference on Software Testing, Verification and
                Validation (ICST), Toronto, ON, Canada, 2024, pp. 317-328.} \\ \textit{doi:}
                \href{https://doi.org/10.1109/ICST60714.2024.00036}{\textit{10.1109/ICST60714.2024.00036}}

            \end{onecolentry}
        \end{samepage}

        \vspace{0.2 cm}

        \begin{samepage}
            \begin{twocolentry}{
                Oct 2023
            }
                \textbf{Precise Lazy Initialization for Programs with Complex Heap Inputs}
            \end{twocolentry}

            \vspace{0.10 cm}
            \begin{onecolentry}
                \mbox{\textbf{\textit{J. M. Copia}}}, \mbox{F. Molina}, \mbox{N.
                Aguirre}, \mbox{M. Frias}, \mbox{A. Gorla}, \mbox{P. Ponzio}. \\
                \textit{2023 IEEE 34th International Symposium on Software Reliability
                Engineering (ISSRE), Florence, Italy, 2023, pp. 752-762.} \\ \textit{doi:}
                \href{https://doi.org/10.1109/ISSRE59848.2023.00080}{\textit{10.1109/ISSRE59848.2023.00080}}

            \end{onecolentry}

        \end{samepage}


        % Juan Manuel Copia, Pablo Ponzio, Nazareno Aguirre, Alessandra Gorla,
        % and Marcelo Frias. 2023. LISSA: Lazy Initialization with Specialized
        % Solver Aid. In Proceedings of the 37th IEEE/ACM International
        % Conference on Automated Software Engineering (ASE '22). Association
        % for Computing Machinery, New York, NY, USA, Article 67, 1–12.
        % https://doi.org/10.1145/3551349.3556965

        \vspace{0.2 cm}

        \begin{samepage}
            \begin{twocolentry}{
                Oct 2022
            }
                \textbf{LISSA: Lazy Initialization with Specialized Solver Aid}
            \end{twocolentry}

            \vspace{0.10 cm}
            \begin{onecolentry}
                \mbox{\textbf{\textit{J. M. Copia}}}, \mbox{P. Ponzio}, \mbox{N.
                Aguirre}, \mbox{A. Gorla}, \mbox{M. Frias}. \\
                \textit{37th IEEE/ACM International Conference on Automated Software
                Engineering (ASE '22). Association for Computing Machinery, New
                York, NY, USA, Article 67, 1–12.} \\ \textit{doi:}
                \href{https://doi.org/10.1145/3551349.3556965}{\textit{10.1145/3551349.3556965}}

            \end{onecolentry}

        \end{samepage}

        %M. Fazzini et al., "Use of Test Doubles in Android Testing: An In-Depth
        %Investigation," 2022 IEEE/ACM 44th International Conference on Software
        %Engineering (ICSE), Pittsburgh, PA, USA, 2022, pp. 2266-2278, doi:
        %10.1145/3510003.3510175.

        \vspace{0.2 cm}

        \begin{samepage}
            \begin{twocolentry}{
                May 2022
            }
                \textbf{Use of Test Doubles in Android Testing: An In-Depth Investigation}
            \end{twocolentry}

            \vspace{0.10 cm}
            \begin{onecolentry}
                \mbox{M. Fazzini}, \mbox{C. Choi}, \mbox{\textbf{\textit{J. M.
                Copia}}}, \mbox{G. Lee}, \mbox{Y. Kakehi}, \mbox{A. Gorla},
                \mbox{A. Orso}. \\
                \textit{2022 IEEE/ACM 44th International Conference on Software
                Engineering (ICSE), Pittsburgh, PA, USA, 2022, pp. 2266-2278.} \\
                \textit{doi:} \href{https://doi.org/10.1145/3510003.3510175}{\textit{10.1145/3510003.3510175}}

            \end{onecolentry}

        \end{samepage}


% =================================================================================================
% -------------------------------------------------------------------------------------------------
% -----------------------------=====>   RESEARCH PROTOTYPES   <=====-------------------------------
% -------------------------------------------------------------------------------------------------
% =================================================================================================

\vspace{0.1 cm}

    \section{Research Prototypes}

    \vspace{0.1 cm}

    \begin{onecolentry}
        \begin{highlightsforbulletentries}

        \item \href{https://github.com/JuanmaCopia/express/}{\textbf{Express}}:
        A tool that automatically infers method preconditions and class
        invariants using search-based algorithms. It specializes in generating
        executable predicates (Java code) as preconditions for complex
        heap-allocated objects. \\
        \vspace{0.1 cm}
        Available at: \href{https://github.com/JuanmaCopia/express/}{github.com/JuanmaCopia/express}

        \vspace{0.2 cm}

        \item \href{https://github.com/facumolina/fixcheck/}{\textbf{FixCheck}}:
        A tool for improving patch correctness analysis in Java. It integrates
        static analysis, random testing, and large language models (LLMs) to
        automatically generate tests that highlight and explain potential patch
        inconsistencies. \\
        \vspace{0.1 cm}
        Available at: \href{https://github.com/facumolina/fixcheck/}{github.com/facumolina/fixcheck}

        \vspace{0.2 cm}

        \item \href{https://github.com/JuanmaCopia/spf-pli/}{\textbf{PLI}}:
        A Symbolic Execution framework built on top of Java Symbolic PathFinder
        (JPF). It enhances symbolic analysis for programs with complex
        heap-allocated inputs and intricate preconditions. PLI addresses key
        challenges in lazy initialization, enabling more precise and efficient
        symbolic reasoning. \\
        \vspace{0.1 cm}
        Available at: \href{https://github.com/JuanmaCopia/spf-pli/}{github.com/JuanmaCopia/spf-pli}

        \vspace{0.2 cm}

        \item \href{https://github.com/JuanmaCopia/SymSolve/}{\textbf{SymSolve}}:
        An efficient bounded exhaustive solver for symbolic structures with
        complex representation invariants. It determines the satisfiability of
        partially symbolic heaps with respect to structural constraints by
        leveraging operational specifications (e.g., repOk routines). \\
        \vspace{0.1 cm}
        Available at: \href{https://github.com/JuanmaCopia/SymSolve/}{github.com/JuanmaCopia/SymSolve}

        \vspace{0.2 cm}

        \item \href{https://github.com/JuanmaCopia/PySEAT/}{\textbf{PySEAT}}:
        A Symbolic Execution engine for Python programs that implements lazy
        initialization to represent heap-allocated objects symbolically. It
        automates test generation for Python codebases, achieving high coverage
        and uncovering potential errors in complex systems. \\
        \vspace{0.1 cm}
        Available at: \href{https://github.com/JuanmaCopia/PySEAT/}{github.com/JuanmaCopia/PySEAT}

        \end{highlightsforbulletentries}
    \end{onecolentry}


% =================================================================================================
% -------------------------------------------------------------------------------------------------
% --------------------------------=====>   PUBLIC TALKS   <=====-----------------------------------
% -------------------------------------------------------------------------------------------------
% =================================================================================================

\vspace{0.1 cm}

    \section{Public Talks}

    \vspace{0.1 cm}


    \begin{twocolentry}{
        Apr 2024
    }
        \textbf{Precise Lazy Initialization for Programs with Complex Heap Inputs}
    \end{twocolentry}

    \vspace{0.1 cm}
    \begin{onecolentry}
        \textit{Workshop,} \href{https://conf.researchr.org/details/icse-2024/klee-2024-papers/23/Precise-Lazy-Initialization-for-Programs-with-Complex-Heap-Inputs}{\textsc{KLEE Workshop on Symbolic Execution}}, \textit{Lisbon, Portugal.}
    \end{onecolentry}


    \vspace{0.2 cm}


    \begin{twocolentry}{
        Oct 2023
    }
        \textbf{Precise Lazy Initialization for Programs with Complex Heap Inputs}
    \end{twocolentry}

    \vspace{0.1 cm}
    \begin{onecolentry}
        \textit{Research Track,}
        \href{https://issre.github.io/2023/program_research.html#research_16}{\textsc{ISSRE 2023}},
        \textit{Florence, Italy.}
    \end{onecolentry}


    \vspace{0.2 cm}


    \begin{twocolentry}{
        Oct 2022
    }
        \textbf{LISSA: Lazy Initialization with Specialized Solver Aid}
    \end{twocolentry}

    \vspace{0.1 cm}
    \begin{onecolentry}
        \textit{Research Track,}
        \href{https://conf.researchr.org/details/ase-2022/ase-2022-research-papers/103/LISSA-Lazy-Initialization-with-Specialized-Solver-Aid}{\textsc{ASE 2022}},
        \textit{Oakland Center, Michigan, USA.}
    \end{onecolentry}


    \vspace{0.2 cm}


    \begin{twocolentry}{
        Sep 2022
    }
        \textbf{LISSA: Lazy Initialization with Specialized Solver Aid}
    \end{twocolentry}

    \vspace{0.1 cm}
    \begin{onecolentry}
        \textit{Oral communication,}
        \href{https://software.imdea.org/es/events/software-seminars/2022/09-27/}{\textsc{IMDEA Software S3 Seminar Series}},
        \textit{Madrid, Spain.}
    \end{onecolentry}


    \vspace{0.2 cm}


    \begin{twocolentry}{
        Mar 2022
    }
        \textbf{A Satisfiability Solver for Symbolic Structures with Complex Representation Invariants}
    \end{twocolentry}

    \vspace{0.1 cm}
    \begin{onecolentry}
        \textit{Oral communication,}
        \href{https://sites.google.com/view/facas2022/charlas-y-actividades-talks-and-activities?authuser=0}{\textsc{FACAS 2022}},
        \textit{La Falda, Córdoba, Argentina.}
    \end{onecolentry}




% =================================================================================================
% -------------------------------------------------------------------------------------------------
% -----------------------------------=====>   Skills   <=====--------------------------------------
% -------------------------------------------------------------------------------------------------
% =================================================================================================

\vspace{0.1 cm}

\section{Technical Skills}

\vspace{0.1 cm}

\begin{onecolentry}
\textbf{Programming:} Python, Java, TypeScript \\
\textbf{Artificial Intelligence:} OpenAI Agents SDK, MCPs, PyTorch, scikit-learn, NumPy, pandas, Matplotlib, data preprocessing, model evaluation \\
\textbf{AWS:} EC2, Lambda, S3, DynamoDB, SQS, SNS, CloudWatch, IAM, CloudFormation, CDK (Infrastructure as Code)
\textbf{Software Engineering:}  Unit Testing, Integration Testing, Code Review \\
\textbf{Java Frameworks and Libraries:} Dagger, Lombok, JUnit \\
\textbf{Research and Analysis:} Static and dynamic analysis, automatic test input generation, symbolic execution, fuzzing, invariant inference \\
\textbf{Tools:} Git, Docker
\end{onecolentry}


\vspace{0.2 cm}

\begin{onecolentry}
    \textbf{Spoken Languages:} English (fluent), Spanish (native), French (fluent)
\end{onecolentry}


\end{document}